\chapter{Estado del Arte}

El uso de la tecnología HSI para la detección de enfermedades fúngicas ha sido ampliamente explorado durante la última década, demostrando gran eficacia en la identificación de hongos y aflatoxinas en maíz, cacahuetes y otros frutos secos \cite{park_hyperspectral_2015, cruz-carrasco_detection_2024}. En el contexto específico de los higos, investigaciones recientes han aplicado técnicas de aprendizaje profundo, tales como las Redes Neuronales Convolucionales (CNN) y Redes Neuronales Totalmente Conectadas (FCNN) \cite{cruz-carrasco_detection_2024}. Estos estudios han establecido precedentes importantes, reportando altas tasas de precisión en la clasificación de píxeles infectados por \textit{Aspergillus flavus} \cite{cruz-carrasco_detection_2024}.

Sin embargo, el análisis reflexivo de estos enfoques revela una importante limitación tecnológica y operativa: la mayoría de estas arquitecturas neuronales profundas utilizan el espectro electromagnético completo \cite{sawant_survey_2020}. Esta dependencia exige un hardware sumamente costoso e impone demandas de alto poder computacional, dificultando su escalabilidad y aplicación práctica en entornos agrícolas reales \cite{sawant_survey_2020}. 

Para contrarrestar la sobrecarga de datos, la literatura sugiere que la selección de bandas espectrales es un paso ineludible \cite{sawant_survey_2020}. Tradicionalmente se han empleado métodos estadísticos como el Análisis de Componentes Principales (PCA) para reducir los datos \cite{chu_classifying_2020}. No obstante, la tendencia actual indica que las técnicas basadas en algoritmos evolutivos están ganando terreno rápidamente debido a su capacidad superior para encontrar soluciones óptimas globales dentro de espacios de búsqueda altamente complejos \cite{sawant_survey_2020}. Este trabajo cubre directamente este vacío, proponiendo el uso de metaheurísticas no solo para procesar la tendencia central del espectro, sino para incorporar el análisis textural (variabilidad), optimizando la detección sin la costosa dependencia de la redundancia espectral total.
